% ============================================
% Assistant Instructional Professor Cover Letter – Harris School of Public Policy (Data Science)
% ============================================

\documentclass[11pt,letterpaper]{letter}

\usepackage{geometry}
\usepackage{microtype}
\usepackage[hidelinks]{hyperref}

\geometry{left=25mm,right=25mm,top=25mm,bottom=25mm}

\signature{Decory Edwards}
\address{
Department of Economics \\
Johns Hopkins University \\
Baltimore, MD 21218 \\
\texttt{dedwar65@jh.edu}
}

\begin{document}

\begin{letter}{Search Committee \\
Harris School of Public Policy \\
University of Chicago \\
Chicago, IL}

\opening{Dear Members of the Search Committee,}

I am writing to apply for the Assistant Instructional Professor position in Data Science at the Harris School of Public Policy. I am a Ph.D.\ candidate in Economics at Johns Hopkins University, advised by Professors Christopher D.~Carroll, Nicholas W.~Papageorge, and Stelios Fourakis, and I expect to receive my degree in May 2026. My research fields are macroeconomics, wealth and income dynamics, and computational economics.

My research agenda focuses on understanding the determinants of wealth inequality and the mechanisms through which heterogeneity in returns to wealth shapes aggregate outcomes. In my job-market paper, I estimate the degree of heterogeneity in individual portfolio returns using micro-data from the Survey of Consumer Finances and simulate a structural model in which households face idiosyncratic return risk. I show that allowing for persistent heterogeneity in returns substantially improves the model’s ability to match the empirical distribution of wealth, offering a tractable way to connect micro-level return dispersion to macroeconomic inequality.

In a second project, I use the Health and Retirement Study to examine the relationship between interpersonal trust and portfolio performance. Building on the literature that finds a hump-shaped relationship between trust and income, I show that a similar relationship holds for individual returns to wealth measured in the HRS data, highlighting a behavioral channel through which social attitudes shape saving and investment outcomes. This work also provides new estimates of the persistent component of returns to net worth, which motivate the calibration of heterogeneity in my structural model.

The Harris School is an excellent fit for my interests in teaching data science in a policy-oriented setting. My research and training rely heavily on data construction, statistical programming, simulation methods, and the translation of complex quantitative results into clear, policy-relevant insights. These skills align closely with Harris’s emphasis on data science, programming, machine learning, and civic technology for social impact.

\textbf{Teaching.} I have experience teaching and supporting courses that integrate data analysis, programming, and applied economic reasoning. My teaching emphasizes reproducible workflows, transparent data practices, and the interpretation of empirical results for policy audiences. At Harris, I would be well prepared to teach courses such as \textit{Introduction to Data and Programming for Public Policy}, \textit{Data Visualization for Policy Analysis}, and applied courses that introduce machine learning tools in a policy context. I place particular emphasis on helping students move from raw data to substantive conclusions while understanding the limitations and assumptions embedded in each step of the analysis.

I am especially drawn to Harris’s focus on using data and technology to address real-world policy problems. I would be enthusiastic about contributing to the School’s instructional mission and about working with students who seek to apply data science tools to questions of governance, regulation, and social welfare.

Thank you very much for your time and consideration. I would be happy to provide any additional information.

\closing{Sincerely,}

\end{letter}
\end{document}
